\section{Bayesian Optimization and Adaptive Stopping}
\label{sec:background}


In black-box optimization, the goal is to find an approximate optimum of an unknown objective function $f : X \to \mathbb{R}$ using a limited number of function evaluations at points $x_1, \ldots, x_T \in X$ where $X$ is the search space and $T$ is a given search budget, potentially chosen adaptively. 
The convention measures performance in terms of \emph{expected simple regret} \citep[Sec. 10.1]{garnett2023bayesian}, given by
\begin{equation}
\label{eqn:bayesopt_problem}
\mathcal{R} = \mathbb{E} \left[\min_{1\leq t\leq T} f(x_t) - \inf_{x\in X} f(x) \right],
\end{equation}
where the expectation is taken over all sources of randomness, including the sequence of points $x_1, \ldots, x_T$ selected by the algorithm.
Bayesian optimization approaches this problem by building a probabilistic model of $f$---typically a Gaussian process (GP)~\citep{rasmussen2006gaussian}, conditioned on the observed data points $(x_t, y_t)_{t=1}^T$, where $y_t = f(x_t)$. 
For each iteration $t=1,\ldots,T$, an acquisition function $\alpha_{t}: X \to \mathbb{R}$ then guides the selection of new samples by carefully balancing the exploration-exploitation tradeoff arising from uncertainty about $f$. 


\subsection{Cost-aware Bayesian Optimization}
\label{sec:background:acquisition}
\emph{Cost-aware Bayesian optimization} \citep{lee2020cost,astudillo2021multi} extends the above setup to account for the fact that evaluation costs can vary across the search space.
For instance, in the hyperparameter tuning example of \Cref{sec:intro}, costs can vary according to the time needed to train a machine learning model under a given set of hyperparameters $x$.
Cost-aware Bayesian optimization handles this by introducing a cost function $c: X \to \mathbb{R}^+$, which may be known or unknown ahead.
The cost function, or observed costs, are then used to construct the acquisition function $\alpha_t$.

In this work, we focus primarily on the \emph{cost-per-sample} formulation \citep{chick2012sequential,cashore2016multi,xie2024cost} of cost-aware Bayesian optimization, which seeks methods with stopping time $\tau$, not necessarily fixed, that achieve a low \emph{expected cost-adjusted simple regret}
\begin{equation}
\mathcal{R}_c = \mathbb{E} \Bigg[\underbrace{\vphantom{\sum_{t=1}^{\tau}}\min_{1\le t \le \tau} f(x_t) - \inf_{x \in X} f(x)}_{\text{simple regret}} + \underbrace{\sum_{t=1}^{\tau} c(x_t)
    }_{\text{cumulative cost}}\Bigg] \label{eqn:cost_per_sample}
    .
\end{equation}
One can also work with the \emph{expected budget-constrained} formulation \citep{xie2024cost}, which incorporates budget constraints explicitly, and seeks algorithms which achieve a low expected simple regret under an expected evaluation budget.
Here, performance is evaluated in terms of
\begin{align}
\mathcal{R} &= \mathbb{E} \left[\min_{1\leq t\leq {\tau}} f(x_t) - \inf_{x\in X} f(x)\right] \nonumber \\
&\text{where } x_1,\ldots,x_{\tau} \text{ satisfy \quad } 
\mathbb{E}\sum\limits_{t=1}^{\tau} c(x_t) \leq B. \label{eqn:cost_aware_bayesopt_problem}
\end{align}
The stopping rules we study can be applied in this setting as well: we discuss this in \Cref{sec:method:theory}.

For both settings, we work with a few acquisition functions---chiefly, the \emph{Pandora's Box Gittins Index (PBGI)} \citep{xie2024cost} and \emph{log expected improvement per cost (LogEIPC)} \citep{xie2024cost} (the log variant \citep{ament2023unexpected} of the expected improvement per cost \citep{snoek2012practical}).
At iteration $t$, let $x_{1:t}:=\{x_1, \dots, x_t\}$ be the evaluated points, $y_{1:t}:=\{y_1, \dots, y_t\}$ be their observed values, $y^*_{1:t}=\min_{1\leq s \leq t} y_s$ be the current best observed value, and $f \mid x_{1:t}, y_{1:t}$ be the posterior distribution of $f$. Then the PBGI acquisition function is defined as
\begin{align}
\label{eqn:PBGI}
&\alpha_t^{\mathrm{PBGI}}(x) := \nonumber\\
&\begin{cases}
g \quad \text{s.t. } % g \text{ solves} \quad
\mathrm{EI}_{f \mid x_{1:t}, y_{1:t}}(x; g) =  c(x), & x\notin x_{1:t}, \\
f(x), & x\in x_{1:t},
\end{cases}
\end{align}
where the value $g$ is the threshold such that the expected improvement over $g$ equals the cost\footnote{For an evaluated point $x$, we set $g = f(x)$ since the posterior at $x$ collapses to a point mass at the observed value $f(x)$ and its evaluation cost can be treated as $0$.} and $\mathrm{EI}_\psi(x; y) = \mathbb{E} \left[\max(y - \psi(x),0)\right]$ is the expected improvement at point $x$ with respect to some random function $\psi : X \to \mathbb{R}$ and a baseline value~$y$.
The LogEIPC acquisition function is defined as
\begin{align}
\label{eqn:EIC} \alpha_{t}^{\mathrm{LogEIPC}}(x) := \log\frac{\mathrm{EI}_{f \mid x_{1:t}, y_{1:t}}(x; y^*_{1:t})}{c(x)},
\end{align}
i.e., the logarithm of the expected improvement divided by the cost.
We also consider classical cost-unaware acquisition functions---\emph{lower confidence bound (LCB)} and \emph{Thompson sampling (TS)}; see \citet{garnett2023bayesian} for details.


\subsection{Adaptive Stopping Rules} \label{sec:background:stopping_rule_review}
%To the best of our knowledge, stopping rules for Bayesian optimization have only been studied in the classical setting without explicit evaluation costs.
To the best of our knowledge, stopping rules for Bayesian optimization typically do not incorporate cost explicitly into the stopping criterion.
We broadly categorize existing methods as follows:

\textbf{Criteria based on convergence or significance of improvement.} 
This includes empirical convergence, namely stopping when the best value remains unchanged for a fixed number of iterations, or the \emph{global stopping strategy (GSS)} \citep{bakshy2018ae}, which stops when the improvement is no longer statistically significant relative to the inter-quartile range (i.e., the range between the 25th percentile and the 75th percentile) of prior observations.
In the multi-fidelity setting, \citet{FourmaniB25} proposed stopping when the high-fidelity surrogate’s estimated optimum has stabilized over a window of iterations, which is related to cost-awareness but differs from our setting with explicitly varying evaluation costs.

\textbf{Acquisition-based criteria.} This includes stopping rules based on acquisition functions such as PI, EI, and KG. These approaches typically stop when the acquisition value falls below a threshold—such as a predetermined constant \citep{lorenz2015stopping,locatelli1997bayesian}, the median of the initial acquisition values \citep{ishibashi2023stopping}, or the cost of sampling \cite{nguyen2017regret,zhou2024corrected,frazier2008knowledge}. They are typically defined via fixed thresholds and are primarily designed for the uniform-cost setting, without explicitly accounting for varying evaluation costs.

\textbf{Regret-based criteria.} This includes stopping rules based on confidence bounds such as \emph{UCB-LCB} \citep{makarova2022automatic} and the \emph{gap of expected minimum simple regrets} \citep{ishibashi2023stopping}. These stop when certain estimated regret bounds fall below a preset or data-driven threshold. The related \emph{probabilistic regret bound (PRB)} stopping rule \citep{wilson2024stopping} stops when estimated simple regret is below a small threshold $\epsilon$ with confidence $1 - \delta$.
In contrast to these forward-looking rules, \citet{li2023not} propose a backward-looking stopping rule that detects when the search has entered a locally convex region and stops once an estimate of local regret falls below a predetermined threshold.

In settings beyond Bayesian optimization, \citet{chick2012sequential} have proposed a cost-aware stopping rules for finite-domain sequential sampling problems with independent values. 
Although this formulation allows for varying costs, it does not extend to general Bayesian optimization settings which use correlated GP models.
